% ════════════════════════════════════════════════════════════
% SÉANCE 7 — ESCALIER
% ════════════════════════════════════════════════════════════
\seance{7}{Escalier}

\subsection{Présentation de l'escalier étudié}

Dans cette étude, l'escalier est considéré comme un \textbf{élément indépendant} de la structure principale. Il constitue un élément essentiel de \textbf{sécurité} et un élément d'\textbf{évacuation} en cas d'incendie. Le dimensionnement portera sur l'escalier de l'étage courant.

\figureplaceholder{escalier-vue-plan}{Vue en plan de l'escalier --- volées et palier intermédiaire}

\figureplaceholder{escalier-coupe}{Coupe en élévation de l'escalier --- paillasse, contremarche, giron}

\paragraph{Caractéristiques :}
\begin{itemize}
  \item Hauteur d'étage : $h_t = 3{,}50$~m
  \item Hauteur à franchir par volée : $h_\text{volée} = 1{,}75$~m
  \item Épaisseur de la paillasse : $e = 12$~cm (valeur retenue selon le critère de flexibilité)
\end{itemize}

\subsection{Pré-dimensionnement par la règle de Blondel}

Le critère de dimensionnement des éléments constitutifs de l'escalier est basé sur la \textbf{règle de Blondel}, donnée par la relation :
\[
  0{,}59 \leq g + 2h \leq 0{,}66~\text{m}
\]

\paragraph{Calculs retenus.} En posant l'hypothèse $\tan\alpha = h/g = 0{,}6$, on obtient :
\[
  g + 2h = 0{,}66~\text{m} \quad\Longrightarrow\quad g + 2(0{,}6\,g) = 0{,}66
\]
\[
  2{,}2\,g = 0{,}66 \quad\Longrightarrow\quad g = 0{,}30~\text{m} \quad ; \quad h = 0{,}18~\text{m}
\]

\paragraph{Détermination du nombre de marches $n$ :}
\[
  n = \frac{h_\text{volée}}{h} = \frac{1{,}75}{0{,}18} = 9{,}72
\]

On retient $n = 10$ marches par volée.

\paragraph{Choix : $n = 10$ marches par volée.}
\[
  h_\text{réel} = \frac{1{,}75}{10} = 0{,}175~\text{m} = 17{,}5~\text{cm}
  \qquad ; \qquad
  g_\text{réel} = 30~\text{cm}
\]

Ainsi, $e = 12$~cm et $a = 10$~cm (largeur du nez de marche).

\paragraph{Longueur de la volée :}
\[
  L = n \cdot g + 2 \cdot \frac{a}{2}
    = 10 \times 30 + 2 \times \frac{10}{2}
    = 310~\text{cm}
    = \mathbf{3{,}10~\text{m}}
\]

\subsection{Descente de charges}

\subsubsection{Charges permanentes (G)}

\paragraph{Poids propre du béton projeté.} L'épaisseur moyenne du béton projeté de la paillasse est donnée par :
\[
  E = \frac{e}{\cos\alpha} + \frac{h}{2}
\]

Avec $\alpha = 30^\circ$ (correspondant à $\tan\alpha \approx 0{,}6$, arrondi pour les calculs) :
\[
  E = \frac{12}{\cos 30^\circ} + \frac{17{,}5}{2}
    = 13{,}856 + 8{,}75
    = \mathbf{22{,}606~\text{cm}}
\]

\[
  g_\text{paillasse} = 0{,}22606 \times 25 = \mathbf{5{,}65~\text{kN/m}^2}
\]

\paragraph{Poids propre du revêtement (marbre).} L'épaisseur équivalente du marbre est :
\[
  e_\text{eq} = \frac{e_m \cdot g + e_m \cdot h}{g}
\]

avec $e_m = 0{,}03$~m (épaisseur du marbre), $g = 0{,}30$~m, $h = 0{,}175$~m :
\[
  e_\text{eq} = \frac{0{,}03 \times 0{,}30 + 0{,}03 \times 0{,}175}{0{,}30}
              = \frac{0{,}009 + 0{,}00525}{0{,}30}
              = \mathbf{0{,}0475~\text{m}}
\]
\[
  g'_\text{marbre} = \gamma_\text{marbre} \cdot e_\text{eq}
                   = 28 \times 0{,}0475
                   = \mathbf{1{,}33~\text{kN/m}^2/\text{ml}}
\]

\paragraph{Total charges permanentes :}
\[
  g_t = g_\text{paillasse} + g'_\text{marbre}
       = 5{,}65 + 1{,}33
       = \mathbf{6{,}98~\text{kN/m}^2/\text{ml}}
\]

\subsubsection{Charges d'exploitation (Q)}
\[
  q = 2{,}5~\text{kN/m}^2
\]

\subsubsection{Combinaison à l'ELU}
\[
  P_u = 1{,}35\,g_t + 1{,}5\,q
      = 1{,}35 \times 6{,}98 + 1{,}5 \times 2{,}5
      = \mathbf{13{,}2~\text{kN/m}^2/\text{ml}}
\]

\figureplaceholder{escalier-schema-statique}{Schéma statique de l'escalier --- charges projetées sur la longueur $L = 3{,}10$~m}

\subsection{Sollicitations}

\paragraph{Moment fléchissant maximal (à mi-portée) :}
\[
  M_u = \frac{P_u \cdot L^2}{8} = \frac{13{,}2 \times 3{,}10^2}{8} = \mathbf{15{,}85~\text{kN}\cdot\text{m/ml}}
\]

\paragraph{Effort tranchant maximal (sur appui) :}
\[
  V_{Ed} = \frac{P_u \cdot L}{2} = \frac{13{,}2 \times 3{,}10}{2} = \mathbf{20{,}46~\text{kN/ml}}
\]

\subsection{Calcul des armatures longitudinales}

Avec $d = 10$~cm, $b = 1$~m, $f_{cd} = 25/1{,}5 = 16{,}67$~MPa, $f_{yd} = 500/1{,}15 = 434{,}78$~MPa.

\paragraph{Moment réduit :}
\[
  \mu = \frac{M_u}{b\,d^2\,f_{cd}}
      = \frac{15{,}85 \times 10^{-3}}{1 \times 0{,}10^2 \times 16{,}67}
      = \mathbf{0{,}0954}
\]

\paragraph{Hauteur de bloc :}
\[
  \alpha = 1{,}25\,(1 - \sqrt{1 - 2\mu})
         = 1{,}25\,(1 - \sqrt{1 - 0{,}1908})
         = \mathbf{0{,}1256}
\]

\paragraph{Bras de levier :}
\[
  z = d\,(1 - 0{,}4\,\alpha)
    = 0{,}10\,(1 - 0{,}4 \times 0{,}1256)
    = \mathbf{0{,}095~\text{m}}
\]

\paragraph{Section d'acier inférieure :}
\[
  A_\text{inf} = \frac{M_u}{z \cdot \sigma_S}
              = \frac{15{,}85 \times 10^{-3}}{0{,}095 \times 434{,}78} \times 10^4
              = \mathbf{3{,}85~\text{cm}^2/\text{ml}}
\]

Pour la largeur de la volée (1,20 m) :
\[
  A_\text{inf} = 3{,}85 \times 1{,}20 = 4{,}62~\text{cm}^2
\]

\encadre{Choix de ferraillage : 6\,HA\,10 ($A_\text{inf,réelle} = 4{,}71$~cm$^2$ pour 1,20~m de largeur)}

\paragraph{Section minimale (Eurocode \S{}9.3.1.1) :}
\[
  A_{s,\min} = 0{,}26 \cdot \frac{f_{ctm}}{f_{yk}} \cdot b_t \cdot d
            = 0{,}26 \times \frac{2{,}6}{500} \times 1 \times 0{,}10 \times 10^4
            = \mathbf{1{,}35~\text{cm}^2/\text{ml}}
\]

La condition $A_\text{inf} \geq A_{s,\min}$ est largement vérifiée.

\paragraph{Armatures de répartition :}
\[
  A_r \approx 1{,}00~\text{cm}^2/\text{ml} \quad\Longrightarrow\quad 4\,\text{HA}\,6/\text{ml}
\]

\figureplaceholder{escalier-ferraillage-paillasse}{Schéma de ferraillage de la paillasse --- armatures principales et armatures de répartition}

\subsection{Ferraillage du béquet}

\paragraph{Caractéristiques du béquet.} L'épaisseur du béquet vaut $e = 7$~cm.

\paragraph{Réaction d'appui à l'ELU :}
\[
  R_u = 2 \cdot V_{Ed} = 2 \times 20{,}5 = \mathbf{41~\text{kN/ml}}
\]

\paragraph{Moment sollicitant :}
\[
  M_u = R_u \cdot \left(l - \frac{a}{2}\right)
\]

avec $l = a + 3~\text{cm} = 13~\text{cm}$ et $a = 10$~cm :
\[
  M_u = 41 \times \left(0{,}13 - \frac{0{,}10}{2}\right) = \mathbf{3{,}28~\text{kN}\cdot\text{m/ml}}
\]

\paragraph{Moment réduit (avec $d = 5$~cm) :}
\[
  \mu = \frac{M_u}{b\,d^2\,f_{cd}}
      = \frac{3{,}28 \times 10^{-3}}{1 \times 0{,}05^2 \times 16{,}67}
      = \mathbf{0{,}078}
\]

\paragraph{Bras de levier :}
\[
  \alpha = 1{,}25\,(1 - \sqrt{1 - 0{,}156}) = \mathbf{0{,}1026}
\]
\[
  z = 0{,}05\,(1 - 0{,}4 \times 0{,}1026) = \mathbf{0{,}0479~\text{m}}
\]

\paragraph{Section d'acier inférieure :}
\[
  A_S = \frac{M_u}{z \cdot \sigma_S}
      = \frac{3{,}28 \times 10^{-3}}{0{,}0479 \times 434{,}78} \times 10^4
      = \mathbf{1{,}57~\text{cm}^2/\text{ml}}
\]

\paragraph{Espacement des armatures :}
\[
  s_t \leq \min\{2h ;\, 20~\text{cm}\}
       = \min\{14~\text{cm} ;\, 20~\text{cm}\}
       = 14~\text{cm}
\]

Ainsi, $s_t = 14$~cm.

\encadre{Choix de ferraillage du béquet : 4\,HA\,6 / ml}

\figureplaceholder{escalier-bequet-ferraillage}{Schéma de ferraillage du béquet --- armatures principales et armatures de répartition}

